\section{Overview: One Call Through the Pipeline}
\label{sec:overview}

\begin{figure}[t]
\centering
\begin{minipage}[c]{0.27\textwidth}
\scriptsize
\begin{verbatim}
fun bump(n) {
  return n + 1;
}

fun main() {
  a = bump(5);
  b = bump(4);
  __voblint_check(a == 6);
  __voblint_check(b == 5);
}
\end{verbatim}
\end{minipage}%
\hfill
\begin{minipage}[c]{0.71\textwidth}
\centering
% Values: voblint --analysis interval --context entry-state --globals warrow
%         --dot docs/readme-figures/contexts.vimp
\resizebox{\linewidth}{!}{%
\begin{tikzpicture}[
  font=\scriptsize,
  pp/.style={draw, rounded corners=1pt, inner sep=2pt, minimum height=3.4ex, minimum width=27mm, align=left},
  bd/.style={pp, fill=black!6, minimum width=24mm},
  flow/.style={-{Stealth[length=1.5mm]}},
  call/.style={-{Stealth[length=1.6mm]}, thick, violet},
  ret/.style={-{Stealth[length=1.6mm]}, dashed, blue!70!black},
  lab/.style={font=\tiny, fill=white, inner sep=1pt}]
  \node[bd] (em) at (0, 0) {entry\_main};
  \node[pp] (p2) at (0, -0.9) {pp2\quad $a, b \in \top$};
  \node[pp] (p3) at (0, -2.3) {pp3\quad $a \in [6,6]$};
  \node[pp] (p4) at (0, -3.7) {pp4\quad \texttt{a == 6}: \verdict{PROVED}};
  \node[pp] (p5) at (0, -4.6) {pp5\quad \texttt{b == 5}: \verdict{PROVED}};
  \draw[flow] (em) -- (p2);
  \draw[flow, gray, dashed] (p2) -- node[lab, left=1mm] {continuation} (p3);
  \draw[flow, gray, dashed] (p3) -- node[lab, left=1mm] {continuation} (p4);
  \draw[flow] (p4) -- (p5);
  \begin{scope}[on background layer]
    \node[draw=gray!60, rounded corners, fit=(em)(p5), inner sep=4pt,
      label={[gray, font=\tiny]above:{\texttt{main}, context $c_0$}}] {};
  \end{scope}
  \node[bd] (b1e) at (5.2, -0.6) {entry\_bump\quad $n \in [5,5]$};
  \node[bd] (b1x) at (5.2, -1.7) {exit\_bump\quad $n \in [5,5]$};
  \draw[flow] (b1e) -- node[lab, right=1mm] {\texttt{return n + 1}} (b1x);
  \begin{scope}[on background layer]
    \node[draw=gray!60, rounded corners, fit=(b1e)(b1x), inner sep=4pt,
      label={[gray, font=\tiny]above:{\texttt{bump}, context $c_1 = [5,5]$}}] {};
  \end{scope}
  \node[bd] (b2e) at (5.2, -3.0) {entry\_bump\quad $n \in [4,4]$};
  \node[bd] (b2x) at (5.2, -4.1) {exit\_bump\quad $n \in [4,4]$};
  \draw[flow] (b2e) -- node[lab, right=1mm] {\texttt{return n + 1}} (b2x);
  \begin{scope}[on background layer]
    \node[draw=gray!60, rounded corners, fit=(b2e)(b2x), inner sep=4pt,
      label={[gray, font=\tiny]above:{\texttt{bump}, context $c_2 = [4,4]$}}] {};
  \end{scope}
  \draw[call] (p2.east) -- node[lab, above, pos=0.45] {enter} (b1e.west);
  \draw[ret] (b1x.west) -- node[lab, below, pos=0.55] {combine} (p3.east);
  \draw[call] (p3.east) -- node[lab, above, pos=0.45] {enter} (b2e.west);
  \draw[ret] (b2x.west) -- node[lab, below, pos=0.55] {combine} (p4.east);
\end{tikzpicture}} % chktex 31
\end{minipage}
\caption{The running example. Left: the VIMP program. Right: its graph as the
analyzer solves it with intervals and entry-state contexts. Each frame is one
procedure in one context; \texttt{bump} gets one copy per entry state. Solid
violet arrows are calls (enter), dashed blue arrows are returns (combine), and
dashed gray arrows link a call site to its continuation. Both checks are
\verdict{PROVED}. Without contexts, the single copy of \texttt{bump} joins both
calls and both checks are \verdict{UNKNOWN}. The states shown are analyzer
output using the warrowing update rule for global unknowns
(\cref{sec:equations}). The drawing omits the intermediate node of
\texttt{bump}'s body and the unreachable node after \texttt{return}.}
\label{fig:running}
\end{figure}

We follow the program of \cref{fig:running} through the analyzer, with
intervals and entry-state contexts. The source program is compiled to a
procedure-aware control-flow graph (CFG). Each procedure $p$ has an entry node
$\ctor{FunctionEntry}\;p$ and a result node $\ctor{FunctionResult}\;p$. The
graph $g$ has two kinds of edges. An \emph{intraprocedural} edge
$(u, a, v) \in \mathit{intra}\;g$ carries an action $a$, such as an
assignment or a guard. A \emph{call} edge
$(u, e, \ctor{FunctionEntry}\;p, k) \in \mathit{calls}\;g$ leads from the call
site $u$ into $p$ and names the continuation $k$ at which the caller resumes.
A predicate $\Glob$ tells which variables are global. A proved forward
simulation relates every source execution to a graph execution with the same
store (\cref{sec:theorem}).

From the graph, the analyzer generates inequations over \emph{local unknowns}
$(v, c)$ for a node $v$ and a context $c$, and over \emph{global unknowns}
without a program point. At \texttt{a = bump(5)}, the analysis's \emph{enter}
operation binds \texttt{n} to $[5,5]$, and this entry state is \emph{routed}
to the callee context $c_1 = [5,5]$; the second call routes to
$c_2 = [4,4]$. \Cref{sec:equations} gives the equations of a call. For these
two calls, the analyzer maintains separate context states and proves both
checks.

To interpret the abstract state in $c_1$, we need the concrete stores
associated with that context. Activation traces and a context policy define
this association (\cref{sec:traces}). The trace of the activation
\texttt{bump(5)} records its caller, the policy admits $c_1$ for that call, and
the \emph{activation collecting
semantics} $\ActColl(v, c)$ collects the final stores of the traces that end
at $v$ and carry $c$. The \emph{node collecting semantics} $\NodeColl(v)$
collects them regardless of context. An analyzer report $\mathit{res}$ defines
two sets at each node $v$. Its report set $\Report_{\mathit{res}}(v)$ is the
union of the stores represented by its abstract states at $v$ across all
contexts. The verdict set $\Verd_{\mathit{res}}(v)$ contains the stores in which
every definite verdict at $v$ holds: the check condition is true for
\verdict{PROVED} and false for \verdict{REFUTED}. Neither verdict asserts
reachability (\cref{sec:theorem}). The soundness argument is a chain of
inclusions at every node:
\begin{equation}
\label{eq:chain}
\text{reached stores}
\;\subseteq\; \NodeColl(v)
\;=\; \textstyle\bigcup_c \ActColl(v, c)
\;\subseteq\; \Report_{\mathit{res}}(v)
\;\subseteq\; \Verd_{\mathit{res}}(v).
\end{equation}
Each inclusion may add unreached stores, which costs precision, but none may
lose a reached store. The main end-to-end theorem
(\cref{sec:theorem}) states the outer inclusions, combining all contexts at
each node.
